\DocumentMetadata{
  pdfversion=2.0,
  pdfstandard=ua-2,
  lang=en-US,
  tagging=on,
  tagging-setup={math/setup=mathml-SE,tabsorder=structure}
}
\documentclass[11pt,letterpaper]{article}
\usepackage[margin=0.68in,includefoot]{geometry}
\usepackage{newcomputermodern}
\usepackage{amsmath,amscd,mathtools}
\usepackage{tikz,adjustbox}
\providecommand{\square}{\mdlgwhtsquare}
\usepackage[hidelinks]{hyperref}
\hypersetup{
  pdftitle={Ergodic Theory PhD Exam, May 2012},
  pdfauthor={University of Florida Department of Mathematics},
  pdfsubject={Graduate mathematics examination},
  pdfdisplaydoctitle=true,
  bookmarksopen=true
}
\setcounter{secnumdepth}{0}
\setlength{\parindent}{0pt}
\setlength{\parskip}{0.28em}
\setlength{\emergencystretch}{3em}
\setlength{\leftmargini}{2.15em}
\setlength{\leftmarginii}{2.65em}
\setlength{\leftmarginiii}{3em}
\pagestyle{empty}
\raggedbottom
\allowdisplaybreaks
\NewTaggingSocketPlug{block/list/label}{examproblem}{%
  \tagstructbegin{tag=Lbl}\tagstructend%
  \tagstructbegin{tag=\LBody}%
  \tagstructbegin{tag=\ProblemHeadingTag,title-o={\ProblemHeadingText}}%
  \tagmcbegin{tag=\ProblemHeadingTag,actualtext={\ProblemHeadingText}}%
  #2%
  \tagmcend%
  \tagstructend%
}
\newenvironment{examproblems}[2]{%
  \def\ProblemHeadingTag{#2}%
  \AssignTaggingSocketPlug{block/list/label}{examproblem}%
  \begin{enumerate}%
  \setcounter{enumi}{#1}%
  \renewcommand{\labelenumi}{\textbf{\theenumi.}}%
  \setlength{\topsep}{0.35em}%
  \setlength{\partopsep}{0pt}%
  \setlength{\itemsep}{0.48em}%
  \setlength{\parsep}{0.18em}%
}{\end{enumerate}}
\newenvironment{examsubparts}{%
  \AssignTaggingSocketPlug{block/list/label}{default}%
  \begin{enumerate}%
  \setlength{\topsep}{0.18em}%
  \setlength{\partopsep}{0pt}%
  \setlength{\itemsep}{0.16em}%
  \setlength{\parsep}{0.1em}%
}{\end{enumerate}}
\newenvironment{examinstructions}{%
  \par\begingroup\itshape
}{%
  \par\endgroup\vspace{0.18em}
}
\makeatletter
\renewcommand\subsection{\@startsection{subsection}{2}{0pt}%
  {-1.25ex plus -.25ex minus -.1ex}%
  {0.55ex plus .1ex}%
  {\normalfont\large\bfseries}}
\renewcommand{\@maketitle}{%
  \newpage\null\vskip -1em%
  {\footnotesize\bfseries
    \href{https://www.ufl.edu/}{University of Florida}, \href{https://math.ufl.edu/}{Department of Mathematics}\par}%
  \vskip -0.5em%
  \tagstructbegin{tag=H1,title={Ergodic Theory PhD Exam, May 2012}}%
  \tagpdfparaOff%
  \tagmcbegin{tag=H1}%
    {\LARGE\bfseries\href{https://gma.math.ufl.edu/exams/ergodic-theory-phd/}{Ergodic Theory PhD Exam}, May 2012\par}%
  \tagmcend%
  \tagpdfparaOn%
  \tagstructend%
  \vskip 0.25em%
  \tagmcbegin{artifact}%
    \hrule height 1pt%
  \tagmcend%
  \vskip 1em%
}
\makeatother
\title{Ergodic Theory PhD Exam, May 2012}
\author{}
\date{}
\begin{document}
\maketitle
\thispagestyle{empty}
\begin{examinstructions}
Do all problems.
\end{examinstructions}
\begin{examinstructions}
Notation: • The statement that \((X,\mathcal{B},\mu,f)\) is a mpt (measure preserving transformation) means that~\(X\) is a set with the probability measure~\(\mu\) defined on the~\(\sigma\)-algebra~\(\mathcal{B}\) and~\(f\) is a bijective, bi-measurable map \(f:X\to X\) which preserves the measure~\(\mu\). • The statement that \((X,d,\mathcal{B},\mu,f)\) is a mph (measure preserving homeomorphism) means that \((X,\mathcal{B},\mu,f)\) is a mpt and in addition, \((X,d)\) is a compact metric space, \(f\) is a homeomorphism, and~\(\mathcal{B}\) is the Borel~\(\sigma\)-algebra.
\end{examinstructions}
\begin{examproblems}{0}{H2}
\def\ProblemHeadingText{Problem 1.}
\item
Let \((X,d)\) be a compact metric space and \(f:X\to X\) a homeomorphism.
\begin{examsubparts}
\item[\textnormal{(a)}]
Define the recurrent set \(\Lambda(f)\) and the nonwandering set \(\Omega(f)\) and show they are both compact, invariant sets.
\item[\textnormal{(b)}]
Show that \(\Lambda(f)\subset\Omega(f)\) and give an example where the inclusion is proper.
\item[\textnormal{(c)}]
Define the omega-limit set \(\omega(x,f)\) of a point \(x\in X\) and show it is a compact, invariant set.
\item[\textnormal{(d)}]
Prove or disprove: For every \(x\in X\), that \(\omega(x,f)\subset\Lambda(f)\).
\end{examsubparts}
\def\ProblemHeadingText{Problem 2.}
\item
Let \((X,d)\) be a compact metric space and \(f:X\to X\) a homeomorphism.
\begin{examsubparts}
\item[\textnormal{(a)}]
Show that there is always an~\(f\)-invariant subset \(Z\subset X\) so that~\(f\) restricted to~\(Z\) is minimal.
\item[\textnormal{(b)}]
Show that there is always a measure~\(\mu\) so that \((X,d,\mathcal{B},\mu,f)\) is a mph.
\end{examsubparts}
\def\ProblemHeadingText{Problem 3.}
\item
Assume \((X,d,\mathcal{B},\mu,f)\) is a mph.
\begin{examsubparts}
\item[\textnormal{(a)}]
Define the support, \(\operatorname{supp}(\mu)\), of the measure~\(\mu\) and show it is a compact, invariant set.
\item[\textnormal{(b)}]
Show that almost every point in \(\operatorname{supp}(\mu)\) is recurrent.
\item[\textnormal{(c)}]
Show that every point in \(\operatorname{supp}(\mu)\) is nonwandering.
\item[\textnormal{(d)}]
Show that \(\mu(\Omega(f))=1\).
\item[\textnormal{(e)}]
If \((X,d,\mathcal{B},\mu,f)\) is ergodic, show that~\(f\) is full orbit transitive on \(\operatorname{supp}(\mu)\).
\item[\textnormal{(f)}]
Show that \(\operatorname{supp}(\mu)\subset\Lambda(f)\), where \(\Lambda(f)\) is the recurrent set of~\(f\).
\end{examsubparts}
\def\ProblemHeadingText{Problem 4.}
\item
State carefully with complete hypothesis.
\begin{examsubparts}
\item[\textnormal{(a)}]
Birkhoff’s pointwise ergodic theorem
\item[\textnormal{(b)}]
von Neumann’s \(L^2\)-mean ergodic theorem
\end{examsubparts}
\def\ProblemHeadingText{Problem 5.}
\item
Let
\[
A=\begin{pmatrix}1&1\\1&0\end{pmatrix}.
\]
\begin{examsubparts}
\item[\textnormal{(a)}]
Define the one-sided subshift of finite type \(\Sigma_A^+\) determined by~\(A\) and the shift map \(\sigma:\Sigma_A^+\to\Sigma_A^+\).
\item[\textnormal{(b)}]
Define one of the standard metrics on \(\Sigma_A^+\).
\item[\textnormal{(c)}]
Show that periodic points of~\(\sigma\) are dense in \(\Sigma_A^+\).
\item[\textnormal{(d)}]
Show that \((\Sigma_A^+,\sigma)\) is forward orbit transitive.
\item[\textnormal{(e)}]
Let \(N_k\) be the number of fixed points of \(\sigma^k\) in \(\Sigma_A^+\). Compute explicitly
\[
\lim_{k\to\infty}\frac{\log(N_k)}{k}
\]
 and justify your answer completely.
\end{examsubparts}
\def\ProblemHeadingText{Problem 6.}
\item
Fix a~\(p\) with \(0<p<1\).
\begin{examsubparts}
\item[\textnormal{(a)}]
Define the Bernoulli measure \(\mu_p\) on the two-sided shift on two symbols \(\Sigma_2\).
\item[\textnormal{(b)}]
Show that \((\Sigma_2,\mathcal{B},\sigma,\mu_p)\) is strong mixing (you don’t have to prove it is a mpt, you may assume that is true).
\item[\textnormal{(c)}]
Define \(\alpha:\Sigma_2\to\mathbb{R}\) by
\[
\alpha(\ldots s_{-1}s_0s_1s_2\ldots)=2s_{-1}-s_0+s_2^2.
\]
 Compute the following limit for \(\mu_p\)-almost every sequence \(\underline{s}\in\Sigma_2\):
\[
\lim_{N\to\infty}\frac{1}{N}\sum_{i=0}^{N-1}\alpha(\sigma^i(\underline{s})).
\]
 and justify your answer completely.
\end{examsubparts}
\def\ProblemHeadingText{Problem 7.}
\item
Let \(f:\mathbb{R}^2\to\mathbb{R}^2\) be defined by \(f(x,y)=(3x,(1/2)y)\).
\begin{examsubparts}
\item[\textnormal{(a)}]
Show that the origin is an unstable fixed point for~\(f\).
\item[\textnormal{(b)}]
Find an invariant, Borel probability measure for~\(f\) (Hint: don’t try too hard).
\item[\textnormal{(c)}]
Show that your invariant measure in (b) is the only invariant, Borel probability measure.
\end{examsubparts}
\end{examproblems}
\end{document}
